%! AP Statistics — Unit 4, Lesson 4.3 — Lesson Plan
\documentclass[10pt]{article}
\usepackage{apstats-article}
\usepackage{apstats-boxes}
\usepackage{graphicx}
\graphicspath{{images/}}

\newcommand{\UnitNumberName}{Unit 4: Probability, Random Variables, and Probability Distributions \quad}
\newcommand{\LessonNumberName}{Lesson 4.3: Introduction to Probability}

\begin{document}
\begin{center}
    {\Large\bfseries \CourseName: \SchoolYear \\
    {\small\bfseries \UnitNumberName \LessonNumberName}}
\end{center}

\begin{tcolorbox}[colback=sky, colframe=navy, boxrule=0.9pt, arc=2mm,
    left=2mm, right=2mm, top=1.5mm, bottom=1.5mm]
    \textbf{Primary Objective:} Students will calculate and interpret probabilities for
    simple events and their complements using sample spaces and the classical definition
    of probability (Big Idea: VAR-4).
\end{tcolorbox}

\renewcommand{\arraystretch}{1.6}
\begin{skillbox}[Priority Ideas \& Skills]{goldbox}
\begin{minipage}[t]{0.48\linewidth}
    \textbf{AP Skills 3.A and 4.B}
    \begin{itemize}
        \item Calculate $P(E)$ for equally likely outcomes (VAR-4.A, Skill 3.A).
        \item Apply the complement rule: $P(E^c) = 1 - P(E)$ (VAR-4.A.4).
        \item Interpret probability as long-run relative frequency (VAR-4.B.1, Skill 4.B).
        \item List sample spaces using tables or tree diagrams (VAR-4.A.1).
    \end{itemize}
\end{minipage}
\hfill
\begin{minipage}[t]{0.48\linewidth}
    \textbf{Key Understandings}
    \begin{itemize}
        \item The sample space is the complete set of non-overlapping outcomes.
        \item Classical probability: $P(E) = \dfrac{\text{outcomes in }E}{\text{total outcomes}}$
              (when all outcomes are equally likely).
        \item Probability is always between 0 and 1, inclusive.
        \item The complement of $E$ contains all outcomes NOT in $E$; $P(E^c) = 1 - P(E)$.
        \item Probability means: in the long run, $E$ occurs in this fraction of trials.
    \end{itemize}
\end{minipage}
\end{skillbox}

\begin{skillbox}[{Vocabulary, Concepts, \& Theorems}]{greenbox}
\begin{minipage}[t]{0.48\linewidth}
    \begin{itemize}
        \item \textbf{Sample space ($S$)} --- the set of all possible non-overlapping outcomes
        \item \textbf{Event} --- a subset of the sample space
        \item \textbf{Equally likely outcomes} --- each outcome has the same probability
        \item \textbf{Classical probability} ---
              $P(E) = \dfrac{|E|}{|S|}$ when outcomes are equally likely
    \end{itemize}
\end{minipage}
\hfill
\begin{minipage}[t]{0.48\linewidth}
    \begin{itemize}
        \item \textbf{Complement ($E^c$ or $E'$)} --- the event that $E$ does \emph{not} occur
        \item \textbf{Complement rule} --- $P(E^c) = 1 - P(E)$
        \item \textbf{Long-run frequency interpretation} --- $P(E)$ = the relative frequency
              with which $E$ occurs in many independent repetitions
        \item \textbf{Tree diagram} --- a tool for listing multi-stage sample spaces
    \end{itemize}
\end{minipage}
\end{skillbox}

\begin{skillbox}[Activate Prior Knowledge \& Spiral Review]{sky}
\begin{minipage}[t]{0.48\linewidth}
    \textbf{Prior Knowledge Check (3--4 min)}
    \begin{itemize}
        \item Review: ``What did our simulation estimate? What is the \emph{true} probability
              the simulation approximates?'' (Lesson 4.2)
        \item Connect: today we calculate the true (theoretical) probability using sample
              spaces and counting, rather than running trials.
        \item Prompt: \textit{``If I roll a fair die, how likely is it to get a number
              greater than 4?''}
    \end{itemize}
\end{minipage}
\hfill
\begin{minipage}[t]{0.48\linewidth}
    \textbf{Connections Forward}
    \begin{itemize}
        \item Two-way tables and conditional probability $\to$ Lessons 4.4--4.5
        \item Mutually exclusive and independent events $\to$ Lesson 4.6
        \item Discrete random variables $\to$ Lessons 4.7--4.8
        \item Long-run frequency interpretation underpins all of Unit 6 inference
    \end{itemize}
\end{minipage}
\end{skillbox}

\begin{skillbox}[Hook \& Lesson]{sky}
\begin{minipage}[t]{0.48\linewidth}
    \textbf{Hook (5--7 min)}\\
    Show a bag with 3 red and 5 blue poker chips. Ask: ``If I reach in without looking,
    what's the chance I pull out a red chip?'' Take guesses. Then ask:
    ``What if I could list every possible draw? Could I calculate the exact probability
    without needing to actually pull chips?''\\
    Transition: \textit{``That's what the classical definition of probability lets us do.''}
\end{minipage}
\hfill
\begin{minipage}[t]{0.48\linewidth}
    \textbf{Lesson Flow (15--18 min)}
    \begin{enumerate}
        \item Define sample space; practice listing $S$ for: one coin, two coins, one die.
        \item Classical probability formula; worked example with a fair die.
        \item Properties: $0 \le P(E) \le 1$; $P(S) = 1$.
        \item Complement rule: what does ``not $E$'' look like in the sample space? How does
              $P(E^c) = 1 - P(E)$ follow from $P(S) = 1$?
        \item Interpretation: ``If we repeated this process many times, how often
              would we expect $E$?'' --- connect to Lesson 4.2 and the LLN.
    \end{enumerate}
\end{minipage}
\end{skillbox}

\begin{skillbox}[Explicit Instruction \& Active Monitoring]{sky}
\begin{minipage}[t]{0.48\linewidth}
    \textbf{Direct Instruction (10 min)}
    \begin{itemize}
        \item Use a two-way table for rolling two dice. Build the $6 \times 6 = 36$ outcome
              table together. Calculate: $P(\text{sum} = 7)$, $P(\text{sum} \geq 10)$,
              $P(\text{NOT sum} = 7)$ using the complement rule.
        \item Emphasize VAR-4.B.1: ``$P(\text{sum} = 7) = 6/36 = 1/6$ means that if we
              rolled two dice over and over, about 1 in every 6 rolls would sum to 7.''
    \end{itemize}
\end{minipage}
\hfill
\begin{minipage}[t]{0.48\linewidth}
    \textbf{Active Monitoring}
    \begin{itemize}
        \item Watch for: students writing probabilities greater than 1 --- review that
              $P(E) \leq 1$ always.
        \item Watch for: not listing the complete sample space (missing outcomes).
        \item Cold-call: ``How do you know your sample space is complete?''
        \item Cold-call: ``In plain English, what does $P(\text{rain}) = 0.3$ mean?''
    \end{itemize}
\end{minipage}
\end{skillbox}

\begin{skillbox}[Group Work \& Differentiation]{redbox}
\begin{minipage}[t]{0.48\linewidth}
    \textbf{Partner Activity (8--10 min)}\\
    Three-tier activity using a standard deck of cards (or a spinner with labeled
    sections). Students list sample spaces, calculate probabilities, use the
    complement rule, and interpret results in context.
\end{minipage}
\hfill
\begin{minipage}[t]{0.48\linewidth}
    \textbf{Differentiation}
    \begin{itemize}
        \item \textit{Support (Tier R)}: Partially completed sample space; students
              fill in and count.
        \item \textit{On-grade (Tier A)}: Build full sample space and calculate several
              probabilities including complement.
        \item \textit{Extension (Tier E)}: Multi-stage sample spaces with tree diagrams;
              interpret probability in context.
        \item \textit{ELL}: Visual sample-space table with color coding.
    \end{itemize}
\end{minipage}
\end{skillbox}

\begin{skillbox}[Individual Work \& Assessment]{redbox}
\begin{minipage}[t]{0.48\linewidth}
    \textbf{Exit Ticket (5 min)}\\
    A spinner is divided into 8 equal sections numbered 1--8.
    \begin{enumerate}
        \item List the sample space.
        \item Calculate $P(\text{even number})$.
        \item Calculate $P(\text{not even number})$ using the complement rule.
        \item Interpret $P(\text{even}) = 0.5$ in a sentence using the phrase
              ``in the long run.''
    \end{enumerate}
\end{minipage}
\hfill
\begin{minipage}[t]{0.48\linewidth}
    \textbf{AP-Style Check}\\
    \textit{``A bag holds 4 green, 3 yellow, and 1 purple marble. If one marble is
    drawn at random, what is the probability that it is NOT yellow?''}
    \begin{enumerate}[label=(\Alph*)]
        \item $3/8$
        \item $1/8$
        \item $5/8$
        \item $4/8$
    \end{enumerate}
    Answer: (C) --- $P(\text{not yellow}) = 1 - 3/8 = 5/8$.
\end{minipage}
\end{skillbox}

\begin{skillbox}[Reinforcement \& Extension]{goldbox}
\begin{minipage}[t]{0.48\linewidth}
    \textbf{Homework / Practice}
    \begin{itemize}
        \item For 4--5 scenarios (dice, cards, spinners, health statistics), list or
              describe the sample space, calculate probabilities, apply the complement
              rule, and interpret results in context.
        \item One AP-style free-response problem requiring full communication of reasoning.
    \end{itemize}
\end{minipage}
\hfill
\begin{minipage}[t]{0.48\linewidth}
    \textbf{Extension / Preview}
    \begin{itemize}
        \item \textit{Extension}: The Addition Rule for mutually exclusive events
              ($P(A \cup B) = P(A) + P(B)$) --- what happens when events are NOT mutually
              exclusive?
        \item \textit{Preview}: Lesson 4.4 introduces two-way frequency tables and
              conditional probability. How might ``given that'' change a probability?
    \end{itemize}
\end{minipage}
\end{skillbox}

\end{document}
